% \abstract{Imbalanced learning presents a significant challenge in machine learning, particularly in applications such as medical diagnosis, fraud detection, and rare event forecasting. Traditional techniques like oversampling and cost-sensitive learning often suffer from computational inefficiency and performance trade-offs. In this study, we propose lightweight data augmentation techniques that balance computational efficiency and model performance. By integrating Mixup augmentation with Deferred Re-weighting (DRW) schedules and evaluating AutoAugment, RandAugment, and Cutout, we systematically analyze their effectiveness on imbalanced datasets. Additionally, we investigate two sampling strategies—Weighted Random Batch Sampler (WRBS) and Weighted Fixed Batch Sampler (WFBS)—to enhance representation learning for minority classes. Experimental results on imbalanced CIFAR-10 and CIFAR-100 demonstrate that our proposed methods significantly improve classification accuracy and generalization while maintaining computational efficiency. Our findings provide a scalable and robust solution for handling imbalanced data, particularly in resource-constrained environments.}

\abstract{Imbalanced learning, when some classes have far fewer examples than others, creates real challenges in areas like medical diagnosis, fraud detection, and rare-event forecasting. Advanced generative methods such as GANs, diffusion models, and DeepSMOTE, can effectively augment minority-class examples but typically demand substantial compute resources and complex pipelines, which limits their practicality in environments with constrained resources. In this work, we focus instead on lightweight data-augmentation techniques that balance speed and accuracy. By combining simple augmentations (Mixup, Cutout, and MAMix) with careful resampling and class-weighting, we achieve significant gains in minority-class performance without extra computational cost. To further strengthen under-represented classes during training, we introduce two sampling strategies, Weighted Random Batch Sampler (WRBS), which draws examples in proportion to class frequency each iteration, and Weighted Fixed Batch Sampler (WFBS), which maintains a constant minority-to-majority ratio in every batch. Evaluated on imbalanced versions of CIFAR-10, CINIC-10, SVHN, and CIFAR-100, our approach consistently improves both minority-class accuracy and overall generalization, while keeping training fast and resource-friendly. These results offer a scalable, robust, and practical solution for handling class imbalance in constrained environments.}


% \keywords{Deep Learning, Imbalanced Classification, Mixup, DRW, Data
% Augmentation}


